% Folder 79, batch 15 — archivists' pages 281–296.
% Pass: Opus 5.5 (claude-opus-5-5), 2026-10-03 — first pass, unchecked against the pages by a human.
%
% Skipped (no \page): 281 (blank sheet); 287, 289, 291, 296 (university
% administrative notices and seminar announcements, typed, with nothing in
% his hand).
% Notation fixed for the batch: A_0, A_\lambda the adjacency relations on the
% set S of vertices; Et(s) (Et_{A_0}(s), Et_0(s)) the star of s; the
% "plus with a dot, indexed by a vertex" written \underset{s}{\dotplus};
% \mathcal{U}, \mathcal{U}_0 … unipotent subgroups; \nu, \mu(\lambda) as on the page.
\documentclass[11pt,a4paper]{article}
\input{../preamble/grothendieck}

\folder{79}
\batch{15}
\pages{281}{296}
\foldertitle{2-polyèdres réguliers triangulés, et modules libres de rangs 2 sur les anneaux $\Lambda=\mathbb{Z}/n\mathbb{Z}$ et $\Lambda=\mathbb{Z}$ : notes manuscrites (s.d.).}
\dating{[à partir de 1980]}
\watermark{Édition de démonstration}

\begin{document}

\section{Étoiles, relations d'adjacence $A_\lambda$ et action de $\lambda$}

\note{Feuillets de brouillon, sans pagination de l'auteur, où les mêmes formules reviennent d'une page à l'autre ; les notations ($S$, $A_0$, $A_\lambda$, $\mathrm{Et}(s)$, $u_\lambda$, $\nu$) sont fixées avant le lot, et l'argument se poursuit au-delà de la page 296.}

\page{282}
\drawing{deux arcs de courbe se faisant face. À gauche, un point $s$ relié à un point $t$ situé sur un arc d'extrémités $t'$ et $t''$, avec à l'extrême gauche un point $\sigma t$ ; à droite, l'image par $\lambda$ : un point $\lambda t$ sur un arc d'extrémités $t'$ et $t''$ (traits renforcés), relié à $\lambda^{-1}s = s'$. Entre les deux, une flèche $\xrightarrow[u_\lambda]{\lambda}$ et une flèche de retour $\xleftarrow[u_{\lambda^{-1}}]{\lambda^{-1}}$.}

$(s,t)\in A_0 \Longrightarrow (s,\lambda t)\in A_\lambda$

$\mathrm{Et}_{A_0}(s) = \mathrm{Et}_{A_\lambda}(\lambda t)$

$t' = \lambda t + \nu$ / $t'' = \lambda t - \nu$

\marginal{$R_0 \subset S\times S^{*}$}
\note{écrit le long du bord droit ; l'exposant de $S$ est peut-être une étoile.}

$\sigma(\lambda t) = \lambda\,\sigma(t)$, \uncertain{lié à} $t$

$(s,t)\in A_0 \Longleftrightarrow (s,\lambda t)\in A_\lambda$

$\forall a$, $\exists f$ triangle, tel que $a\in f$

\struck{\ill{}} $\operatorname{card} A = N/2$

$\operatorname{card} F = (N/2)\,|$\note{le dernier signe est un trait vertical, peut-être un début de fraction laissé en suspens.}

$(a,f)$ un nb $\nu(N/2)$

donc il y a $\nu(N/2)/3$ faces

Si $\nu=1$, ça ferait $N/6$ faces

le stabilisateur d'une face serait d'ordre $6$

\note{trait de séparation sur toute la largeur.}

\drawing{triangle de sommets $\lambda s$, $\lambda t$, $\lambda u$.}

$\lambda t \pm \nu$ est incident à $t$, avec, sous le signe, $s' = \lambda^{-1}s$

$\lambda(t \underset{s}{\dotplus} \zeta) \pm \nu$ est incident : $t \underset{s}{\dotplus} \zeta$
\[
  \lambda t \underset{s'}{\dotplus} \lambda^{2}\zeta \pm \nu
\]
\[
  \zeta' = \lambda^{2}\zeta \pm \nu, \qquad \zeta = \lambda^{-2}\zeta' \pm \lambda^{-2}\nu
\]

$\lambda^{-2}\nu = \nu^{-1}$\note{en bas à gauche, isolé.}

\marginal{\uncertain{irréelle} \struck{dans} au beau milieu de délire mathématico-philoso\supplied{phique}}
\note{à l'encre bleue, le long du bord droit, au stylo et non au crayon comme les calculs ; fragment de phrase sans lien visible avec eux, sans début ni fin sur la page.}

\page{283}
\drawing{suite de petites figures : un triangle curviligne de base $s\,t$ ; un triangle subdivisé ; un triangle de sommets $s$, $t$, $u$ ; un segment $s\,t$ prolongé vers deux points marqués d'une flèche, avec $t'$ et $s'$ ; plusieurs segments issus de $s$ et $t$ superposés.}

$\forall s\in S$, $\exists s'\in S$ tel que $\mathrm{Et}_{A_0}(s) = \mathrm{Et}_{A}(s')$

\drawing{une ligne $s$ — $t$ — $\lambda t$ — $s' = \lambda^{-1}s$, avec un point $t'$ au-dessus de $\lambda t$ (qui surcharge un \struck{$\lambda t$} d'abord mal placé) ; puis une ligne $s$ — $t$ — $t'$, et un point $s'$ isolé.}

\page{284}
$ga\mathcal{U} = g'a\mathcal{U}$ ssi $g\mathcal{U} = g'\mathcal{U}$ i.e. $g'^{-1}g\in\mathcal{U}$

$a^{-1}(g'^{-1}g)a$ \struck{\ill{}} $\in\mathcal{U}$ ssi $g'^{-1}g\in\mathcal{U}$

$a^{-1}\mathcal{U}a$

\drawing{plusieurs étoiles de triangles (un hexagone subdivisé autour d'un sommet $t$, d'autres sommets marqués $s$, $s_-$, $t_-$), un carré coupé par une diagonale, un quadrilatère $s$, $t$, $s_-$, $t_-$ ; des arêtes renforcées.}

$x \quad \lambda x$

$G_{s_0} = G_s = G_{g.s_0}$, avec sous $G_{s_0}$ : $\mathcal{B}$, et sous $G_{g.s_0}$ : $\operatorname{int}(g)\mathcal{B}$\note{les deux égalités verticales sont écrites $\parallel$ ; l'indice du premier $G$ est surchargé.}

$x\wedge y = e \Longrightarrow x\wedge \underbrace{\lambda y}_{y'} = \lambda e$

$\lambda^{-1}x\wedge y' = e$

\struck{$\lambda x\wedge y$}

$\mathrm{Et}_{A_0}(s) \xrightarrow{\ \lambda\ } \mathrm{Et}_{A_0}(\lambda^{-1}s)$\note{le $s$ du second membre surcharge un autre signe.}
\[
  \mathrm{Et}_{A_0}(s) = \{\,y \mid x\wedge y = e\,\}, \qquad
  \mathrm{Et}_{A_0}(\lambda^{-1}s) = \{\,y' \mid \lambda^{-1}x\wedge y' = e \text{ i.e. } x\wedge y' = \lambda e\,\}
\]
\[
  y \longmapsto y' = \lambda y
\]
\[
  y \mapsto y + x, \qquad y' \longmapsto y' + \lambda^{-1}x
\]

\struck{$\lambda x$} $y_1 = y + x$ \qquad \struck{$y'_1 = y' + \lambda$}

$\rho^{-1}(u)\,\rho^{-1}(a^{-1})\,a^{-1}$
\[
  \underbrace{\lambda y_1}_{y'_1} = \underbrace{\lambda y}_{y'} + \lambda x =
\]
\[
  y'_1 = y' + \lambda^{2}(\lambda^{-1}x)
\]

\drawing{un tétraèdre, un losange coupé par sa diagonale, quelques arcs.}

$\rho_0\mathcal{U} = \rho_\infty\mathcal{U}$\note{l'indice du premier $\rho$ est surchargé.} \qquad $\tau_0\tau_1$

$\rho_\infty^{-1}\tau_0\in\mathcal{U}$ \qquad $\tau_1$

$\tau_0\underbrace{\tau_1\tau_0}_{\rho_\infty}\in\mathcal{U}$ \qquad $\mathrm{Et}(s) = \langle\tau_1,\tau_\infty\rangle(s,t)$\note{lecture de $\tau_\infty$ incertaine.}

$\tau_0\tau_1\tau_0$ \qquad $s,t$

$\mathrm{Et}_{A_0}(t) = \mathrm{Et}_{A_\lambda}(\lambda t)$

$\{s,t\}\in A_0 \Longleftrightarrow \{s,\lambda t\}\in A_\lambda$\note{le $t$ du premier membre surcharge un signe biffé.}

$\mathrm{Et}(s)\cap\mathrm{Et}(\lambda t)$

$\{s,\lambda^{-1}u\}\in A_0$ \qquad $\{t,\lambda^{-1}u\}\in A_0$

$(s,t)\in A_\lambda$ \qquad $\{s,\lambda u\}\in A_\lambda$ \qquad $(t,u)\in A_\lambda$

\page{285}
$a\rho^{-1}a\rho\in\widetilde{\mathcal{U}}\rho^{-1}(\widetilde{\mathcal{U}})$ \qquad $\rho(a)a\in\rho(\widetilde{\mathcal{U}})\cdot\widetilde{\mathcal{U}}$

\struck{$a\rho^{-1}(a)\in$}

$au^{-1}\rho^{-1}a\rho v\in\widetilde{\mathcal{U}}\rho^{-1}(\widetilde{\mathcal{U}})$

\struck{\ill{}} $a^{-1}\rho^{-1}a\rho v\in\widetilde{\mathcal{U}}\,(\rho^{-1})(\widetilde{\mathcal{U}})$

\struck{$\widetilde{\mathcal{U}}$} $v^{-1}\rho^{-1}\widetilde{\mathcal{U}}\rho v$\note{relié par une flèche à la ligne précédente.}

$\rho(ua)a\in\mathcal{U}$

\struck{$\rho(a)a^{-1}(u)$} $\rho(ua)(av)$

$\rho(a)av\in\mathcal{U}_1\mathcal{U}_0$

$\rho$ \qquad $\mathcal{U}_0\ \mathcal{U}_1\ \mathcal{U}_\infty$\note{un arc relie $\mathcal{U}_0$ à $\mathcal{U}_1$.}

$\varepsilon_0 = \sigma_\infty^{-1}\rho$ \qquad
$\sigma_\infty(\varepsilon_0) = \rho\,\sigma_\infty$, \quad
$\rho_\infty(\varepsilon_0) = \rho\,\sigma_\infty^{-1}$ \quad définis par $\omega$\note{l'indice de $\rho_\infty$ surcharge un autre signe.}

$\rho : e_1\longmapsto -e_2$, \quad \struck{$e_\tau$} : $e_2\longmapsto e_1+e_2$ \qquad $\omega$

$\begin{pmatrix}\lambda & \mu\\ 0 & \lambda^{-1}\end{pmatrix}$

\marginal{$\mathrm{Et}_{A_0}(s)\cap\mathrm{Et}_{A_0}(\lambda^{-1}t)\neq\emptyset$ ; \quad $s$ — $s'$ — $\lambda^{-1}t$ ; \quad $s,\lambda^{-1}t$ ; \quad \struck{$\lambda^{-1}s = s$} ; \quad \struck{$\lambda^{-1}s = t$}}
\note{écrit perpendiculairement aux calculs, le long du bord droit.}

\drawing{un carré en pointillés avec ses deux axes et deux diagonales fléchés (les vecteurs $\pm e_1$, $\pm e_2$, $\pm(e_1+e_2)$), étiquetés $s$, $t$, $\lambda^{-1}s$, $\lambda^{-1}t$ ; un hexagone avec trois diagonales formant un triangle inscrit ; un carré ; deux fois $E_0$.}
\[
  \begin{pmatrix}1 & 0\\ \mu & 1\end{pmatrix}
  \begin{pmatrix}1 & \mu'\\ 0 & 1\end{pmatrix}
  = \begin{pmatrix}1 & \mu'\\ \mu & 1+\mu\mu'\end{pmatrix}
\]
\[
  \rho = \begin{pmatrix}0 & 1\\ -1 & 1\end{pmatrix}
\]
\[
  \underbrace{\begin{pmatrix}0 & 1\\ -1 & 1\end{pmatrix}
  \begin{pmatrix}\lambda & 0\\ 0 & \lambda'\end{pmatrix}}_{\begin{pmatrix}0 & \lambda'\\ -\lambda & \lambda'\end{pmatrix}}
  \underbrace{\begin{pmatrix}1 & -1\\ 1 & 0\end{pmatrix}
  \begin{pmatrix}\lambda & 0\\ 0 & \lambda'\end{pmatrix}}_{\begin{pmatrix}\lambda & -\lambda'\\ \lambda & 0\end{pmatrix}}
  = \begin{pmatrix}\lambda\lambda' & 0\\ \lambda\lambda'-\lambda^{2} & \lambda\lambda'\end{pmatrix}
\]
\note{sous la première accolade le manuscrit porte $\begin{pmatrix}0 & -\lambda'\\ -\lambda & \lambda'\end{pmatrix}$, le premier $-$ peu net ; on a gardé le signe moins dans cette note et non dans la formule faute de le lire sûrement. Le $\lambda'$ du second facteur surcharge un $\lambda$.}
\[
  = \begin{pmatrix}1 & 0\\ 1-\lambda^{2} & 1\end{pmatrix}
  \qquad \mu' = 0, \quad \mu = 1-\lambda^{2}
\]
\note{devant $\mu' = 0$, un signe biffé.}

\drawing{tracé à l'envers par rapport au reste de la page : une ligne $\lambda^{-1}(s)$ — $t$ — $s$, et au-dessus $\rho^{-1}(s)$, avec un sommet d'où partent trois segments.}

\page{286}
$\mu(\lambda)\in(\mathbb{Z}/n\mathbb{Z})^{*}$

$\nu(\lambda)\in(\mathbb{Z}/n\mathbb{Z})/\pm1$ \qquad $\nu$

\begin{tikzcd}
S \arrow[r, "\lambda^{-1}", "\sim"'] \arrow[d, "\mathrm{Et}_0"', "\wr"] & S \arrow[d, "\mathrm{Et}_0", "\wr"'] \\
\mathcal{E} \arrow[r, "\sim", "\lambda_{\mathcal{E}}"'] & \mathcal{E}
\end{tikzcd}

\note{à droite, l'indice de $\mathrm{Et}$ est mal formé ($0$ ou $\lambda$) ; à l'extrémité de la flèche horizontale du bas, un mot biffé \ill{} ; les deux $\mathcal{E}$ surchargent chacun un signe biffé.}

\begin{tikzcd}
S \arrow[r, "\mathrm{Et}_0", "\sim"'] \arrow[dr, "\mathrm{Et}_\lambda"'] & \mathcal{E} \arrow[d, "\wr", "\lambda_{\mathcal{E}}"'] \\
 & \mathcal{E}
\end{tikzcd}

\note{l'étiquette de la diagonale surcharge un signe biffé.}

\drawing{un hexagone dont tous les sommets sont joints entre eux.}

$E$ \qquad \struck{$R_E$}\note{lecture du signe biffé incertaine.}

$\omega_E$ \qquad $\nu$ \qquad $T_\lambda$

$\lambda\,\mathrm{Et}_0(s) = \mathrm{Et}_0(\lambda^{-1}s)$

$\lambda\,\mathrm{Et}_0(s) = \mathrm{Et}_\lambda(s)$

$(s,t)$ \struck{\uncertain{adj}} $A_0$-adj. $\Longrightarrow$ $\lambda t$ $A_\lambda$-adj. à $\lambda t \underset{s,A_0}{\dotplus}\nu \overset{?}{\Longrightarrow}$

\struck{$\nu$} \qquad \struck{$(s,t)$} $\nu$ dépend de \emph{$A_0$} et de $\lambda$

\drawing{une étoile de six triangles autour d'un sommet.}

$E$ \struck{\uncertain{adjacents}} \struck{$(s,t)\in A_0$}

$\omega_{A_0}$ \qquad $\nu$
\[
  \lambda \longmapsto \omega_{A_0}^{\nu}, \qquad
  \lambda' \longmapsto \omega_{A_0}^{\nu'}, \qquad
  \lambda\lambda' \mid
\]
\note{les exposants $\nu$, $\nu'$ sont écrits au-dessus de $\omega_{A_0}$ ; la ligne $\lambda\lambda'$ reste inachevée.}

\struck{$A_{\lambda}$} $T_\lambda(A_0) =$

$\lambda t \dotplus \nu$ $A_\lambda$-adj. à \struck{\ill{}} $\lambda t$

$T_\lambda^{2}(A_0) = \lambda.(A_0)$

\struck{$\lambda$}$^{-1}t -$ \qquad $T_\lambda$

$(s,t)\in A_0 \Longleftrightarrow (\lambda s,\lambda^{-1}t)\in A_0$\note{dans le premier membre, un $\lambda$ biffé devant $s$.}

$\Updownarrow\,?$ \qquad $(\lambda\mu s, \lambda\mu(t))\in A_0$\note{lecture de cette ligne incertaine.}

$\parallel$ \qquad $(\mu(\lambda)s, \mu^{-1}\lambda^{-1}t)\in A_0$, avec sous $\mu^{-1}\lambda^{-1}t$ : $(\lambda\mu)^{-1}t$

\page{288}
\drawing{un arc de cercle.}
\[
  u_\lambda s = \lambda^{-1}s, \qquad u_\lambda t = \lambda t, \qquad
  \begin{pmatrix}\lambda^{-1} & 0\\ 0 & \lambda\end{pmatrix}
\]

$u_\lambda$ normalise $\mathcal{U}_s$ et $\mathcal{U}_t$

\struck{$\lambda$}$t$ incident à $\lambda t \underset{s'}{\pm}\nu$

$u_\lambda t = \lambda t$ incident à $\lambda^{2}t \underset{s}{\pm}\nu$ \qquad $\lambda t$.

$t$ adj. : $\lambda t \underset{s'}{\pm}\nu$\note{ligne entourée.}

$\lambda t$ est $A_\lambda$-adj. : $\lambda t\pm\nu$

$\lambda t$ est $A_0$-adj. $\lambda^{-1}(t\pm\nu) = $ \struck{$\lambda$}$t\pm\lambda^{2}\nu$\note{sous le dernier membre, un mot biffé \ill{} ; l'exposant $\lambda^{-1}$ du premier membre surcharge un autre signe.}

$\lambda t$ adj. :

\marginal{$u_\lambda$}
\note{à gauche, le long d'un grand arc qui embrasse les quatre dernières lignes.}

\page{290}
$N/n$ étoiles ou sommets, $N/2$ arêtes

Si chaque arête par $A_\lambda$ est contenue dans $\alpha$ \struck{étoiles} triangles $(\alpha\geq1)$, i.e. dans \struck{$\alpha$} étoiles, alors le nb de couples $(a,E)$ est
\[
  (N/2)\,\alpha = (N/n)\cdot\beta
\]
où $\beta$ nb d'arêtes dans une étoile donnée
\[
  \beta = \frac{n\alpha}{2}
  \begin{cases}
    \dfrac{n}{2} & \text{si } \alpha=1 \\
    n & \text{si } \alpha=2
  \end{cases}
\]
\note{en regard du premier cas : « $n$ pair, et alors on trouve $n$ diagonales ». Devant $\beta$, un signe biffé.}

\drawing{un pentagone dont tous les sommets sont joints entre eux.}
\[
  \nu(\lambda)\,\nu(\lambda^{-1}) \overset{?}{=} 1
\]
\[
  \lambda^{-2}\nu(\lambda) = \nu(\lambda^{-1}) \overset{?}{=} \nu(\lambda)^{-1}
\]
donc $\nu(\lambda)^{2} = \lambda^{2}$ $\Longrightarrow$ $\boxed{\nu(\lambda)\in(\mathbb{Z}/n\mathbb{Z})^{*}}$\note{le premier $\nu$ de la ligne $\nu(\lambda)\nu(\lambda^{-1})$ surcharge un autre signe.}

\page{292}
$\begin{pmatrix}\lambda & 0\\ 0 & \lambda^{-1}\end{pmatrix}$ normalisateur de $\mathcal{U}_0$

opère sur $\mathcal{U}_0$ par $\boxed{\lambda^{2}}$

et $\lambda : \mathrm{Et}(s)\xrightarrow{\ \sim\ }\mathrm{Et}(\lambda^{-1}s)$ qui transforme $\omega_s$ pour une orientation fixée en $\omega_{\lambda^{-1}s}^{\boxed{\lambda^{2}}}$ i.e.
\[
  \boxed{\lambda.(a\dotplus\zeta) = \lambda a + \lambda^{2}\zeta}
\]

Donc on ne \uncertain{voit ici} que $\lambda^{2}$, pas $\lambda$ ! Mais la relation d'incidence entre $a\underset{s}{\dotplus}\zeta\in\mathrm{Et}(s)$ et $\lambda a\underset{\lambda^{-1}s}{\dotplus}\zeta'$ est
\[
  \zeta' = \lambda^{2}\zeta \pm \lambda
\]
i.e. les quantités $\zeta' - \lambda^{2}\zeta$ forment l'ens $\{\pm\lambda\}$ !\note{$\lambda^{2}$ souligné ; dans $\lambda.(a\dotplus\zeta)$, le $a$ surcharge un autre signe.}

\note{trait de séparation sur toute la largeur.}

$\zeta,\zeta'$ incidents \qquad $\zeta' = \pm\lambda$ \qquad \struck{\ill{}} \qquad $t'\pm\nu_\lambda$

\drawing{un segment $s$ — $t$ — $\lambda t = t'$ — $\lambda^{-1}s = s'$, inscrit dans un grand cercle tracé légèrement ; à droite, un arc et un segment aboutissant à $s'$.}
\[
  \zeta'_1 = \lambda^{2}\zeta + \lambda, \quad \zeta'_2 = \lambda^{2}\zeta - \lambda
  \qquad
  \begin{cases}
    \zeta'_1 + \zeta'_2 = 2\lambda^{2}\zeta \\
    \zeta'_1 - \zeta'_2 = 2\lambda
  \end{cases}
\]

$\mu(\lambda)$

\note{dans un cadre :} $a$ incident : $\lambda(a)\underset{\lambda^{-1}(s)}{\pm}\lambda$, avec sous $a$ : $\Uparrow$ $\mathrm{Et}(s)$, et après $\lambda^{-1}(s)$ un signe \ill{} suivi de $\Lambda$

\page{293}
\[
  \mathrm{Et}(s) \xrightarrow[\ y\,\mapsto\,\lambda y\ ]{\ \lambda\ } \mathrm{Et}(\lambda^{-1}s)
\]
\[
  \mathrm{Et}(s) \simeq \{\,y \mid x\wedge y = e\,\}, \qquad
  \mathrm{Et}(\lambda^{-1}s) = \{\,y' \mid \lambda^{-1}x\wedge y' = e \text{ i.e. } x\wedge y' = \lambda e\,\}
\]

Relation entre $y,y'$ \qquad $\boxed{\lambda(y\dotplus\zeta) = \lambda y + \lambda^{2}\zeta}$ \emph{à prouver}
\[
  y\wedge y'\in\pm e
\]

$\mathrm{Et}(s)$ torseur sous $\Lambda$ par
\[
  y\dotplus\zeta = y + \zeta x
\]

$\mathrm{Et}(\lambda^{-1}s)$ torseur sous $\Lambda$ par
\[
  y'\dotplus\zeta = y + \zeta(\lambda^{-1}x) = y + (\lambda^{-1}\zeta)x
\]
\note{au second membre le manuscrit porte $y$, non $y'$.}

\emph{NB} $\lambda(y\dotplus\zeta) = \lambda(y+\zeta x) = \lambda y + \lambda\zeta x = \lambda y \dotplus (\lambda^{2}\zeta)$

\struck{Relation $y_0$} entre $\mathrm{Et}(s)$, $\mathrm{Et}(\lambda^{-1}s)$
\[
  y\wedge y'\in\pm e
\]
\[
  y = y_0\dotplus\zeta = y_0 + \zeta x, \qquad
  y' = y'_0\dotplus\zeta' = y'_0 + (\lambda^{-1}\zeta')x
  \qquad y'_0 = \lambda y_0 \text{ d'où } y_0\wedge y'_0 = 0
\]
\[
  y\wedge y' = (\lambda\zeta - \lambda^{-1}\zeta')\,e
\]
\[
  \varphi(\zeta,\zeta') = \lambda\zeta - \lambda^{-1}\zeta' \in \pm1
\]
\[
  \begin{cases}
    \varphi(\zeta,\zeta') = \varphi(\zeta+\alpha,\ \zeta'+\lambda^{2}\alpha) \\
    \varphi(\zeta,\zeta')
  \end{cases}
\]
\[
  \boxed{\zeta' = \lambda^{2}\zeta \pm \lambda} \qquad \zeta'
\]

\page{294}
\[
  (y_0 + \zeta x)\wedge(y'_0 + \zeta'\lambda^{-1}x)
  = (\lambda\zeta - \lambda^{-1}\zeta' + \theta)\,e
\]
\note{dans le premier facteur, le $\lambda$ de $\lambda\zeta$ surcharge un autre signe.}

\struck{\ill{}} $\zeta - \zeta' + \theta \in \pm1$
\[
  \boxed{\lambda\zeta - \lambda^{-1}\zeta' + \theta \in \pm1}
\]
\[
  \zeta' = \underbrace{\lambda^{2}\zeta + \lambda\theta}_{\text{dépend des origines choisies}} \pm \lambda
  \qquad \zeta\in\mathbb{Z}/n\mathbb{Z}
\]

ops $\theta = 0$\note{lu « ops », sans doute pour « on peut supposer ».}
\[
  \boxed{\mathrm{Et}(s) \xrightarrow[\ \sim\ ]{\ \lambda\ } \mathrm{Et}(\lambda^{-1}s)}
  \qquad \lambda\zeta - \lambda^{-1}\zeta'
\]

$(s,\lambda t)\in A_0 \Longleftrightarrow \{\lambda t,\lambda^{-1}s\}\in A_0$\note{le $\lambda t$ du premier membre surcharge un signe biffé.}
\[
  \lambda\zeta - \lambda^{-1}\zeta' \in \pm1
  \qquad \alpha+\zeta \quad \lambda\alpha + \zeta' \quad \zeta - \lambda^{-1}\zeta'
\]

\struck{$\lambda(\zeta+\alpha)$} \qquad $t$ \qquad $\lambda t$

\struck{\ill{}} $\lambda(t+\zeta) =$ \struck{$\lambda(t+\zeta)$} $\lambda t + \lambda^{2}\zeta$

\struck{\ill{}} $t+\zeta$ \qquad $\lambda t + \zeta'$

$\lambda\zeta$ \struck{$=$} $- \lambda^{-1}\zeta'\in\pm1$
\[
  \lambda(\alpha+\zeta) - \lambda^{-1}(\lambda^{2}\alpha + \zeta')\in\pm1
\]
\[
  (t+\alpha)+\zeta = t + (\alpha+\zeta), \qquad
  \lambda(t+\alpha) + \zeta' = \lambda t + (\lambda^{2}\alpha + \zeta')
\]
\note{les deux dernières lignes en bas à droite ; devant la première, un mot biffé \ill{}.}

\page{295}
\[
  T_\lambda(A_{\lambda'}) = A_{\lambda\lambda'}
\]
\note{l'indice $\lambda'$ du premier membre surcharge un autre signe.}

$\lambda$ \emph{défini} par $\nu$ \qquad
$\mathrm{Et}_0(\lambda^{-1}t) = \{\,s\in S \mid \{s,\lambda^{-1}t\}\in A_0\,\}$

$\lambda'$ \emph{défini} par $\nu'$ \qquad
$\phantom{\mathrm{Et}_0(\lambda^{-1}t)} = \{\,s\in S \mid \{s,t\}\in A_\lambda\,\}$

\note{la première égalité est barrée d'un trait oblique léger.}

$\{s,t\}\in A_\lambda \Longrightarrow \{\lambda$

$\nu(A_0,\lambda)$ défini par la condition

$\{s,t$

$E\in\mathcal{E}$ \qquad $A_{\in E}$\note{lecture de l'indice incertaine.}

\struck{\ill{}}

\note{le bas de la page porte l'étiquette imprimée du courrier interne de l'université, sans texte de sa main. Les notes du lot s'arrêtent ici ; la suite éventuelle est au-delà de la page 296.}

\end{document}
